\chapter{Conclusions and Future Work}
\label{chap:conclusions}

This thesis presented a comprehensive investigation into the application of the $A^*$ search algorithm for the optimization of Multi-Layer Perceptrons. By transitioning from a discretized state-space exploration to structured neighbor generation, we established a framework that challenges the traditional reliance on gradient-based methods for specialized neural architectures.

\section{Summary of Findings}

The experimental campaign, spanning from preliminary hyperparameter tuning to large-scale benchmarks, has yielded several critical insights:

\begin{itemize}
    \item \textbf{Optimal Search Configuration:} Preliminary studies in Chapter \ref{chap:hyperparameters_tuning} demonstrated that a vanilla $A^*$ approach is insufficient for deep architectures. The introduction of \textbf{Beam Search} and \textbf{Layer-Wise Kernel Reshaping} proved essential to mitigate the curse of dimensionality, providing a balance between computational feasibility and convergence depth.
    
    \item \textbf{High-Precision Convergence:} Across the Iris Flower and Noisy Sine benchmarks, the $A^*$ framework—specifically the kernel-based variants—consistently achieved training losses several orders of magnitude lower than the Adam baseline. In the Iris task, the framework reached absolute zero training loss, demonstrating a level of parameter refinement that gradient-based plateaus often prevent.
    
    \item \textbf{The Crossover Law of Generalization:} A significant finding of this research is the "crossover" effect observed as network complexity increases. While Adam remains superior in speed, the \textbf{Layer-Wise $A^*$ strategy} began to outperform Adam in test-set accuracy on larger architectures (e.g., Iris Big Net), suggesting that structured search identifies weight configurations with superior generalization margins.
    
    \item \textbf{Robustness and Stability:} Unlike Adam, which exhibited oscillations in non-linear regression tasks, the $A^*$ framework provided a smooth and deterministic-like descent. Even in complex classification tasks like Wine Quality, where Adam achieved deeper training minima, the $A^*$ strategies demonstrated significantly lower variance across independent runs, indicating a lower sensitivity to initial weight distributions.
\end{itemize}



\section{Future Work}

While the results validate the potential of search-based optimization, the following areas remain open for future exploration:

\begin{itemize}
    \item \textbf{Asynchronous and Parallel Search:} The primary bottleneck remains the wall-clock time required for neighbor evaluation. Future work will investigate the implementation of asynchronous $A^*$ search on multi-GPU clusters to parallelize the expansion of the Open Set, potentially reducing training time by orders of magnitude.
    
    \item \textbf{Hybrid Neuro-Evolutionary Optimizers:} There is significant potential in a multi-stage optimizer. A hybrid approach could utilize Adam to rapidly reach the proximity of a local minimum and then transition to a \textbf{Layer-Wise $A^*$ Search} for high-precision refinement of the weights where gradients become vanishingly small.
    
    \item \textbf{Adaptive Quantization and Dynamic Step-Sizing:} Building upon the tuning results in Chapter \ref{chap:hyperparameters_tuning}, future versions of the algorithm could implement a dynamic quantization factor that increases resolution automatically as the loss slope flattens, allowing the search to "zoom in" on global minima without the memory overhead of a globally fine grid.
    
    \item \textbf{Extension to Recurrent and Convolutional Architectures:} This study focused on MLPs. Extending the kernel-based search logic to convolutional kernels or recurrent units would provide insights into whether search-based stability translates to spatial and temporal feature extraction.
\end{itemize}

\section{Final Remarks}

In conclusion, this research demonstrates that $A^*$ search is not merely a theoretical curiosity in neural network training but a viable high-precision optimizer. When properly constrained via Beam Search and structured kernels, it offers a robust alternative to gradient-based methods, particularly in scenarios where training depth, stability, and generalization are prioritized over raw execution speed.